\documentclass[12pt]{article}
\usepackage{fullpage}
\usepackage[T1]{fontenc}
\usepackage[utf8]{inputenc}
\usepackage{lmodern}
\usepackage{microtype}
\usepackage{amsmath,amssymb}
\usepackage{hyperref}
\usepackage{enumitem}

\title{Descriptive Language, Realism, and the Case for Theoreticism}
\author{}
\date{}

\begin{document}
\maketitle

\begin{abstract}
This essay analyses the function of descriptive language regarded as an
instrument for talking about the world, and on that basis evaluates the
thesis that operationalism and instrumentalism should be replaced by
\emph{theoreticism}---the view that inquiry proceeds within a framework of
theories whose aim is genuine explanation, not the mere correlation of
observables. I argue for three connected conclusions, and I grade each
carefully. (i)~The descriptive use of language carries, built into its
very structure, a commitment to a standard of correctness that the act of
describing does not itself secure; this \emph{token-independence} of the
success condition is what distinguishes description from other uses of
language and is the precondition of truth and falsity (\emph{deductively
established}). (ii)~Token-independence does not by itself entail
\emph{recognition-transcendence}, and so does not by itself settle the
dispute between the realist and the sophisticated (Dummettian)
anti-realist who identifies truth with idealized warranted
assertibility. I therefore do \emph{not} claim a conclusive deduction of
semantic realism; I argue instead that semantic realism is the
\emph{best explanation} of the full constellation of features of
descriptive practice, and is in that sense abductively---not
analytically---superior to its sophisticated rival. (iii)~Whether one
accepts the strong (realist) or the weaker (token-independence) reading,
both \emph{refute operationalism and instrumentalism in their strict
forms} and \emph{establish the theory-ladenness limb of theoreticism};
the explanatory-aim limb is supported abductively, via a no-miracles
argument tempered---against the pessimistic meta-induction---by retreat to
\emph{structural} realism. I mark throughout exactly where the argument is
deductive and where abductive, so that the endorsement of theoreticism is
not overstated.
\end{abstract}

\section{Aim, method, and standard of success}

The problem asks me to do four things: analyse the descriptive function
of language; ask whether that function argues for realism; ask whether
such realism justifies replacing operationalism and instrumentalism by
theoreticism; and reach a determinate verdict. Because this is a
philosophical thesis and not a formal theorem, ``proof'' here means a
\emph{valid argument from premises that are either analytic of the
relevant concepts or else more plausible than their denials}, together
with an explicit accounting of where the argument is deductive and where
it is merely abductive. I will accordingly state premises explicitly,
draw the inference, and then test each premise against the strongest
objection I can find. A claim is allowed to stand only if it survives
that test; where it survives only as the best available explanation
rather than as a necessary truth, I say so. I hold myself to this
anti-overclaiming standard strictly: in particular I will \emph{not}
classify as ``analytic'' or ``conclusive'' any step that in fact embeds a
substantive metaphysical commitment a competent opponent rejects.

I fix terminology first, because the entire argument turns on not letting
the key terms drift. Two distinctions introduced here do load-bearing
work later and must not be conflated.

\begin{description}[leftmargin=1.6em,style=nextline]
\item[Descriptive (representative) use of language.] The use of a
declarative sentence to \emph{say how things stand}---to represent a
state of affairs as obtaining---so that what is said is correct if things
are as represented and incorrect otherwise.
\item[Token-independence of a success condition.] A descriptive use $R$
is \emph{token-independent} when whether $R$ succeeds is not constituted
by, and does not logically follow from, the mere act of producing or
accepting \emph{this} token of $R$. (Even a verificationist who equates
truth with idealized assertibility grants token-independence: an
individual assertion can be warrantedly assertible-in-the-limit yet
mistaken now.)
\item[Recognition-transcendence (semantic realism).] $R$'s success
condition is \emph{recognition-transcendent} when it may obtain or fail
to obtain independently of \emph{any} epistemic access we could in
principle have---so that $R$ may be true though forever unverifiable, or
false though ideally assertible. Semantic realism is the thesis that
descriptions have recognition-transcendent success conditions.
\item[Realism.] The conjunction of (a)~\emph{ontological realism}: there
is a world whose constitution does not depend on our beliefs, theories,
or descriptions of it; and (b)~\emph{semantic realism} as just defined.
\item[Sophisticated anti-realism (Dummett).] Truth is not
recognition-transcendent but is \emph{idealized/superassertibility}: $R$
is true iff $R$ is, or would be, warrantedly assertible under the
indefinite extension and refinement of our evidence and would not
thereafter be defeated. This view \emph{accepts} token-independence and
\emph{rejects} recognition-transcendence.
\item[Operationalism.] The doctrine that the meaning of a scientific
concept is \emph{exhausted by} the operations used to measure it: a
concept just \emph{is} the corresponding set of measuring operations.
\item[Instrumentalism.] The doctrine that scientific theories are not
candidates for truth or falsity but only \emph{instruments} (calculating
devices, inference tickets) for systematising and predicting observable
phenomena.
\item[Theoreticism.] The doctrine that all inquiry---including
observation and measurement---is conducted from within a framework of
theories, and that the proper aim of such theories is to \emph{explain},
i.e.\ to say truly what underlies and produces the phenomena, not merely
to correlate observables.
\end{description}

\section{The function of descriptive language as an instrument}

\subsection{The hierarchy of functions of language}

Language performs several distinguishable functions. Following the
analysis associated with B\"uhler and refined by Popper, we may
distinguish at least four:

\begin{enumerate}[label=(F\arabic*),leftmargin=2.4em]
\item\label{F1} the \emph{expressive} (symptomatic) function: a verbal
utterance expresses or symptomises the inner state of the speaker;
\item\label{F2} the \emph{signalling} (release) function: an utterance
stimulates or releases a reaction in a hearer;
\item\label{F3} the \emph{descriptive} (representative) function: an
utterance describes a state of affairs, and is thereby true or false;
\item\label{F4} the \emph{argumentative} function: utterances stand in
relations of support, contradiction, and inference to one another.
\end{enumerate}

I previously asserted a strict inclusion ordering on these functions. I
now state more carefully what is, and what is not, defensible, because
\S4.2 leans on the relation between \ref{F3} and \ref{F4}.

\begin{description}[leftmargin=1.6em,style=nextline]
\item[Lemma 0 (Dependence of contentful inference on description).]
\emph{Material} argumentation in sense \ref{F4}---inference whose premises
and conclusions are evaluated as \emph{correct or incorrect about
something}, and which can be \emph{refuted by how things turn out}---is
possible only over items that are descriptive in sense \ref{F3}.
\end{description}

\noindent\emph{Argument, and the limit of the claim.} One must concede the
critic's counterexample: inference in a \emph{purely uninterpreted
calculus} (the instrumentalist's ``inference-ticket'' picture, in which
sentences are mere marks manipulated by formation and transformation
rules) is argumentation-without-description. So the blanket claim
``\ref{F4} presupposes \ref{F3}'' is false as stated. But uninterpreted
calculi are precisely \emph{not} what scientific theorising is: a calculus
becomes a vehicle of inquiry only when its tokens are
\emph{contradicted}, \emph{confirmed}, or \emph{revised in the light of
how the world is found to be}. Such evaluations---``this consequence is
\emph{wrong}'', ``the data \emph{refute} this''---presuppose that the
tokens \emph{say something} that can be at odds with the facts, i.e.\ that
they are descriptive (\ref{F3}). Hence what \ref{F4} presupposes \ref{F3}
is not formal manipulation but \emph{the truth-directed argumentative
practice that science actually exhibits}. This restricted claim is what
\S4.2 needs, and it is immune to the uninterpreted-calculus
counterexample because that counterexample lacks the very feature
(answerability to the world) that makes a calculus a candidate scientific
theory rather than a game. I therefore do not rely on a strict global
inclusion of the four functions; I rely only on Lemma~0.
$\qquad\blacksquare$

\begin{description}[leftmargin=1.6em,style=nextline]
\item[Lemma 1 (Truth presupposes description).] The predicates ``true''
and ``false'' apply to a use of language only insofar as that use is
descriptive in sense \ref{F3}.
\end{description}

\noindent\emph{Argument.} An utterance is true when things are as it says
they are, and false otherwise (Tarski's adequacy condition: ``snow is
white'' is true iff snow is white). This biconditional has application
only where there is something the utterance \emph{says is the
case}---a represented state of affairs that may or may not obtain. A
purely expressive cry \ref{F1} or a purely releasing signal \ref{F2} says
nothing is the case; there is no represented state of affairs to be
compared with anything; so ``true'' and ``false'' lack purchase. Hence
truth-evaluability and descriptive use stand or fall together. (Note this
lemma is neutral between realist and anti-realist accounts of \emph{what}
truth consists in; it concerns only the bearers of truth-value.)
$\qquad\blacksquare$

\subsection{The internal structure of a description}

To describe is to deploy a representation $R$ that \emph{purports to
correspond} to a state of affairs $S$, in such a way that the description
\emph{succeeds} if $S$ obtains and \emph{fails} if it does not. Three
components are therefore constitutive of the descriptive function:

\begin{enumerate}[label=(\roman*),leftmargin=2.2em]
\item a \emph{representation} $R$ (the content of what is said);
\item a \emph{represented item} $S$ (a state of affairs);
\item a \emph{success condition}: $R$ is correct iff $S$ obtains, where
``$S$ obtains'' is not the same fact as ``$R$ is asserted'' or ``$R$ is
accepted''.
\end{enumerate}

\begin{description}[leftmargin=1.6em,style=nextline]
\item[Lemma 2 (Bipolarity / the possibility of error, scoped).] Any
\emph{contingent} descriptive use of language admits both correctness and
error: it is constitutive of a contingent description $R$ that it could be
\emph{false} and could be \emph{true}, and that which it is is not
constituted by the making of $R$. Equivalently: every contingent
description has a \emph{token-independent} success condition.
\end{description}

\noindent\emph{Argument, and the scope restriction.} Suppose some
contingent use $R$ were descriptive yet could not be false---nothing would
count as $R$'s misrepresenting how things stand. Then there would be no
contrast between $R$'s success condition obtaining and its failing; the
correctness of $R$ would not depend on anything beyond the making of $R$.
But then $R$ would not \emph{say how things contingently stand} as opposed
to merely \emph{occurring}; it would have collapsed into the expressive or
signalling functions \ref{F1}--\ref{F2}, contradicting the assumption that
it is a contingent description. Hence a contingent descriptive use is
\emph{bipolar} and token-independent.

The scope restriction to \emph{contingent} description is forced and I
mark it explicitly, because the unrestricted lemma has clear
counterexamples: ``$2+2=4$'' and ``all bachelors are unmarried'' are
descriptive and truth-apt yet cannot be false. The reductio fails for
them precisely because their truth is \emph{not} a contrast with a
possible falsity grounded in how the contingent world is; their
correctness is fixed by meaning or mathematical structure. This matters
for honesty later: \emph{Lemma~2 may be applied only to contingent
empirical descriptions, not to non-contingent metaphysical or analytic
theses.} In particular, the self-refutation arguments of \S6 against
constructivism must not lean on Lemma~2 applied to a non-contingent
thesis; I reconstruct them accordingly there.
$\qquad\blacksquare$

Lemma~2, so scoped, establishes exactly \emph{token-independence}: the
intelligibility of a contingent description requires a standard of
correctness not guaranteed by the act of describing. It does \emph{not}
yet establish recognition-transcendence. The whole delicacy of \S3 lies
in not confusing these.

\section{From the descriptive function to realism}

\subsection{What the descriptive function establishes deductively}

\begin{enumerate}[label=(P\arabic*),leftmargin=2.4em]
\item\label{P1} If $R$ is a contingent description, then $R$ has a
token-independent success condition (Lemmas~1--2).
\item\label{P2t} A token-independent success condition is one whose
obtaining is not constituted by our asserting or accepting $R$.
\item\label{C0} Therefore contingent descriptive language presupposes
that descriptions answer to a standard of correctness external to the act
of describing them.
\end{enumerate}

The inference \ref{P1}, \ref{P2t} $\Rightarrow$ \ref{C0} is valid and its
premises are analytic of ``contingent description''. \emph{This much is
conclusive.} It already does substantial work in \S4: both operationalism
and strict instrumentalism deny exactly \ref{C0} for the cases that
matter.

\subsection{Why \ref{C0} does not by itself yield semantic realism}

The earlier draft of this essay slid from token-independence (\ref{C0})
to recognition-transcendence (semantic realism) by an allegedly
``analytic'' premise. The critic is right that this equivocates, and I
withdraw the claim of conclusiveness. The sophisticated (Dummettian)
anti-realist \emph{grants} \ref{C0}: she agrees that an individual
assertion can be mistaken, that ``verified now'' and ``true'' come apart.
She identifies truth not with present verification but with
\emph{superassertibility}---warranted assertibility that would survive the
indefinite extension and refinement of evidence and never thereafter be
defeated. On her view the success condition is token-independent
(no one assertion settles it) yet \emph{not} recognition-transcendent (it
cannot outrun \emph{all} possible evidence). \ref{C0} is silent between
this view and realism. So:

\begin{description}[leftmargin=1.6em,style=nextline]
\item[Honest interim verdict.] The descriptive function of language
\emph{conclusively} establishes token-independence (\ref{C0}); it does
\emph{not} conclusively establish semantic realism (recognition-trans\-cen\-dence).
The latter must be argued for as the better of two live positions, not
read off the concept of description.
\end{description}

\subsection{The abductive case for realism over sophisticated anti-realism}

I now argue that realism is nonetheless the \emph{rationally preferable}
position---superior as an explanation of descriptive and scientific
practice---without pretending the argument is a deduction.

\begin{description}[leftmargin=1.6em,style=nextline]
\item[Lemma 3$'$ (The pressure on superassertibility).] The
identification of truth with superassertibility faces three difficulties
that realism avoids, jointly making realism the better explanation of how
descriptive language functions.
\end{description}

\noindent\emph{(i) The intelligibility of evidence-transcendent
contingencies.} We appear to understand, and routinely traffic in,
contingent claims for which we have no guarantee of even ideal evidence:
``there was an even number of dinosaurs at the instant of the impact'',
``some star outside our light cone has a planet''. These are
\emph{bipolar} (Lemma~2) and we treat them as determinately true or
false. The realist explains this trivially: their truth is fixed by how
the (recognition-transcendent) world is. The superassertibility theorist
must either declare them \emph{neither} true nor false (revising classical
logic and our ordinary practice with them) or stretch ``ideal evidence''
until it is indistinguishable from ``how the world is anyway''---at which
point the position has collapsed into realism. Neither horn is costless;
the realist horn is none.

\noindent\emph{(ii) The explanatory role of truth.} Realism explains the
instrumental success of our descriptions: our predictions tend to come
out right \emph{because} the descriptions latch onto how things mind-inde\-pen\-dently
are. Superassertibility, by tying truth to our own idealized verdicts,
makes the match between assertibility and the world either a primitive
coincidence or itself in need of a realist explanation. The
no-miracles intuition (\S5) here favours realism not as proof but as
better explanation.

\noindent\emph{(iii) The non-collapse of the actual/ideal gap.} The
caloric and heliocentrism examples (well-confirmed falsehood; long-unrecognised
truth) show only the \emph{actual}-evidence/truth gap, which the
sophisticated anti-realist grants; I no longer present them as refuting
her. But they motivate the realist reading of the \emph{further}
gap: science repeatedly treats its own best-warranted-so-far theories as
candidates for being false, and seeks deeper theories that may overturn
them \emph{even where the predecessor was empirically impeccable}. This
practice of holding even idealized current warrant hostage to an
independent world is more naturally read as realist commitment than as a
detour through superassertibility.

\noindent None of (i)--(iii) is a deduction; the Dummettian can pay the
prices and remain consistent (her position is not self-refuting---see \S6,
Objection~1, corrected). But the prices are real and cumulative, and no
comparable prices attach to realism. I therefore conclude, \emph{at the
grade of inference to the best explanation}, that semantic realism is
preferable. $\qquad\blacksquare$

\subsection{Verdict on the first sub-question, graded}

\begin{description}[leftmargin=1.6em,style=nextline]
\item[Verdict on realism.] The descriptive use of language \emph{does}
supply an argument bearing on realism, of two distinct strengths.
\emph{(a)~Conclusively}, it establishes token-independence (\ref{C0}): one
cannot coherently make contingent descriptions while denying that they
answer to a standard external to the act of describing. \emph{(b)~Abductively}
(Lemma~3$'$), it supports the stronger semantic realism
(recognition-transcendence) as the best explanation of descriptive and
scientific practice, against the consistent but costly Dummettian rival.
I do not claim (b) as a deduction.
\end{description}

This two-tier verdict is what the \S1 standard demands. It also fixes the
strength available downstream: any step in \S4 that needs only \ref{C0}
inherits conclusive force; any step that needs full semantic realism
inherits only the (high-grade) abductive force of Lemma~3$'$. I track
this dependence explicitly.

\section{From realism to theoreticism}

I argue that the results of \S2--\S3 tell against operationalism and
instrumentalism and in favour of theoreticism. I separate what follows
from token-independence alone (conclusive) from what needs full semantic
realism (abductive).

\subsection{Against operationalism (needs only \ref{C0})}

\begin{description}[leftmargin=1.6em,style=nextline]
\item[Claim 4.] Operationalism is false: scientific concepts are not
exhausted by the operations used to measure them.
\end{description}

\noindent\emph{Argument.} (i)~\emph{Over-multiplication.} If a concept
just is its measuring operations, distinct operations define distinct
concepts; length-by-rod, length-by-triangulation, length-by-light-travel
would be three concepts, not one quantity measured three ways. But science
treats them as measuring \emph{the same} length and asks whether the
methods \emph{agree}---a question presupposing a single quantity, fixed
independently of any one operation, to which the operations are
answerable. The report ``these operations agree'' is bipolar (Lemma~2):
they \emph{might disagree}, and then we conclude one is in error, not that
the world changed its number of lengths. This requires only
token-independence (\ref{C0}): length must be fixed otherwise than by the
operation-token, on pain of the agreement-question's being unintelligible.
No appeal to recognition-transcendence is needed, so this objection is
\emph{conclusive}. (ii)~\emph{Operation presupposes theory.} Specifying a
measuring operation already deploys descriptive, theory-laden concepts
(``rigid rod'', ``simultaneous'', ``undisturbed sample''); the operation
cannot define the concept because understanding the operation already
requires the concept. Operationalism thus either misindividuates concepts
(i) or presupposes the descriptive content it claims to eliminate (ii).
Both are settled by \ref{C0} alone.
$\qquad\blacksquare$

\subsection{Against instrumentalism (strict form: needs only \ref{C0} and Lemma~0)}

\begin{description}[leftmargin=1.6em,style=nextline]
\item[Claim 5.] Strict instrumentalism is untenable: theories cannot be
construed as mere predictive devices with no truth-value.
\end{description}

\noindent\emph{Argument.} (i)~\emph{Observation statements are already
descriptive.} The predictions the instrumentalist grants are
\emph{statements about observables}---``the pointer reads $5$''. By
Lemmas~1--2 these are contingent descriptions: bipolar, truth-evaluable,
token-independent (\ref{C0}). Truth is thus already in play at the
observational level the instrumentalist accepts; she cannot quarantine
truth to observation while denying it to theory, since both are
descriptive in the same sense. (ii)~\emph{Theories are asserted, tested,
refuted.} Scientists assert theories, deduce consequences, regard rivals
as \emph{contradicting} them, and revise them under refutation. By
Lemma~0 (not by any global F1--F4 inclusion), this truth-directed
argumentative practice presupposes that theories \emph{say something} that
can be at odds with the world---i.e.\ that they are descriptive. A device
that can be \emph{contradicted by evidence} is being treated as a
description, not a slide rule. This too needs only \ref{C0} and Lemma~0,
so the refutation of \emph{strict} instrumentalism is \emph{conclusive}.

A residual move: ``Grant predictions are true or false, but a theory is a
mere efficient summary of predictions, with no surplus content to be true
\emph{of}.'' This fails: a theory with no content beyond its actual
predictions could not be \emph{tested by novel} predictions, yet
testability by novel cases is the mark of a theory; the surplus content is
what does the predicting, and (being descriptive) carries its own
token-independent success condition.

I flag explicitly: this refutes \emph{strict} instrumentalism (no
truth-value for theories). It does \emph{not} by itself refute van
Fraassen's constructive empiricism, which grants theories truth-values but
restricts the \emph{aim} of science to empirical adequacy. That is a
distinct target, addressed in \S4.3 and \S6, Objection~2, where only
abductive force is claimed.
$\qquad\blacksquare$

\subsection{For theoreticism}

\begin{description}[leftmargin=1.6em,style=nextline]
\item[Claim 6.] Theoreticism is the position that accommodates \S2--\S3
and survives the objections that sink operationalism and strict
instrumentalism.
\end{description}

Theoreticism asserts: (T1)~inquiry is conducted from within a framework of
theories (no theory-neutral operational or observational ground floor);
and (T2)~the aim of theories is genuine explanation---true description of
the structures producing the phenomena---not mere correlation.

\emph{(T1), suitably stated, is established by \S4.1(ii) and the
theory-ladenness of measurement.} I must, however, correct an overclaim
the critic identified. The strong slogan ``there is no theory-neutral
notion of the observable'' overreaches against van Fraassen, who
distinguishes the theory-ladenness of the \emph{language} we use to report
observations (which he grants) from an \emph{observer-relative}
(anthropocentric, not theory-relative) notion of what is \emph{observable}
by unaided human perception. His observable/unobservable line is drawn by
human physiology, not by theory, and is not refuted merely by noting that
our \emph{descriptions} are theory-laden. So (T1) should be stated as the
defensible claim it is: \emph{the concepts deployed in observation and
measurement, and the warrant we extend to observation reports, are
theory-laden}; there is no \emph{epistemically privileged,
theory-independent} stratum that could serve as operationalism's or
strict instrumentalism's foundation. This is established by \S4.1(ii)
(measurement-concepts are theory-laden) together with the holistic point
that what we count as a good observation (a clean cloud-chamber track, a
genuine parallax) depends on accepted theory. It does \emph{not} require
denying van Fraassen's physiological observable/unobservable distinction,
and I do not so deny it. (T1) so stated is conclusive; the stronger slogan
is withdrawn.

\emph{(T2) is supported abductively, and is the locus of the residual
disagreement with van Fraassen.} Granting realism (\S3.3) or even just
the truth-aptness of theories (\S4.2), the question is whether the
\emph{aim} of science is explanatory truth about structure (T2) or mere
empirical adequacy. To say the world is thus-and-so at the level of
unobserved structure \emph{is} to explain why the observable
correlations hold (they hold because that structure obtains); the
correlations are the observable shadow of the posited structure, and are
predicted \emph{because} it is posited. The case for (T2) is therefore an
inference to the best explanation, taken up in \S5.
$\qquad\blacksquare$

\section{The abductive case for the explanatory aim, and its principled limits}

The argument for (T2) is a version of the no-miracles argument: the
sustained novel-predictive and unifying success of mature theories would
be inexplicable (a ``miracle'') unless those theories were at least
approximately true of the structures they posit; the explanatory-truth
aim best accounts for the actual selection pressures in science (novel
prediction, unification, inter-method agreement). I must, in honesty, set
this against its standard and serious defeater, which the essay itself
half-invites by citing caloric as a verified falsehood.

\begin{description}[leftmargin=1.6em,style=nextline]
\item[The pessimistic meta-induction (Laudan).] The history of science is
a graveyard of once-successful theories now held false---phlogiston,
caloric, the optical and electromagnetic ethers. If past successful
theories were false, the inference from current success to truth is
undercut; so the no-miracles argument over-reaches.
\end{description}

\noindent\emph{Reply: retreat to structural realism.} I do not claim, and
the critic rightly warns against claiming, that mature theories are true
\emph{in their full ontology}. The defensible (T2) is \emph{structural}:
what is retained across the theory-changes Laudan cites is, characteristically,
\emph{mathematical-structural} content---the form of the laws and the
relational structure---even when the posited entities (caloric, ether) are
discarded. Fresnel's equations survive into Maxwell's theory though
Fresnel's ether does not; the inverse-square \emph{structure} of static
fields survives; conservation structures persist. The no-miracles
argument, applied to \emph{structure} rather than to full ontology, both
explains novel success and survives the meta-induction, because it is
exactly the structural content that is preserved across the revolutions.

\begin{description}[leftmargin=1.6em,style=nextline]
\item[Graded verdict on (T2).] The aim of science is best understood as
the discovery of true \emph{structural} descriptions that explain the
phenomena---an explanatory aim, hence theoreticism's (T2), not mere
empirical adequacy. This is \emph{abductive}, and its grade is high but
\emph{tempered}: full-blown entity-realism for every posit is \emph{not}
claimed; the meta-induction forces the structural qualification, and the
empirical adequacy theorist (van Fraassen) remains a consistent, if (I
argue) less explanatory, rival. I therefore retract the earlier
unqualified phrase ``highest available grade'' in favour of: a strong but
defeasible inference to the best explanation, structurally restricted to
survive the pessimistic meta-induction.
\end{description}

\subsection{The strength of the endorsement, stated honestly}

Assembling the grades:
\begin{itemize}[leftmargin=1.6em]
\item Token-independence \ref{C0}: \emph{conclusive}.
\item Refutation of \emph{strict} operationalism (Claim~4) and
\emph{strict} instrumentalism (Claim~5): \emph{conclusive}, since they
need only \ref{C0} and Lemma~0.
\item Theory-ladenness limb (T1, suitably stated): \emph{conclusive};
the over-strong ``no theory-neutral observable'' slogan is withdrawn.
\item Semantic realism (recognition-transcendence): \emph{abductive}
(Lemma~3$'$), preferable to but not a deductive defeat of sophisticated
anti-realism.
\item Explanatory-aim limb (T2): \emph{abductive and structurally
qualified} (\S5).
\end{itemize}

\begin{description}[leftmargin=1.6em,style=nextline]
\item[Final verdict.] The descriptive use of language to talk about the
world \emph{does} supply an argument in favour of realism: conclusively
for token-independence, and abductively (as best explanation) for full
semantic realism. This in turn \emph{does} justify replacing
operationalism and instrumentalism by theoreticism---\emph{deductively}
as regards the refutation of their strict forms and the theory-ladenness
of inquiry (which need only the conclusive results), and \emph{abductively}
as regards the explanatory, structure-describing aim of theory (which
needs the further, best-explanation results). The replacement is
justified; its negative and theory-ladenness components are demonstrated,
and its explanatory-aim component is supported by a structurally tempered
inference to the best explanation. I make no claim of conclusiveness for
the latter, in keeping with the \S1 standard.
\end{description}

\section{Objections and replies}

\paragraph{Objection 1 (global anti-realism / constructivism).} ``The
world we describe is constituted by our conceptual scheme; there is no
scheme-independent success condition.'' \emph{Reply (corrected).} I no
longer run the tu-quoque ``this thesis is itself a bipolar description,
hence self-refuting'', for two reasons the critic identifies: (a)~the
thesis is a \emph{non-contingent, metaphysical} claim, so Lemma~2 (scoped
to contingent description) does not apply to it; (b)~a global anti-realist
can apply her own warranted-assertibility semantics \emph{reflexively}
without contradiction, so the tu-quoque is non-probative. The defensible
reply is the relocation argument, which I now develop rather than merely
assert. Constructivism does not abolish a standard of correctness external
to particular descriptions; it \emph{relocates} it: descriptions now
answer to \emph{the scheme together with the deliverances it organises}.
But (i)~the fact that a given description fits or fails to fit the scheme
is itself something the describer can get \emph{wrong} (we misapply our
own concepts), so token-independence \ref{C0} survives at the level of
scheme-relative correctness; and (ii)~the constructivist owes an account
of what makes one scheme's deliverances recalcitrant where another's are
not---why the world ``pushes back''---and the most economical account
reintroduces a scheme-independent constraint, i.e.\ partial realism. Thus
constructivism does not threaten the conclusive result \ref{C0}, and it
inherits, against the realist, exactly the abductive contest of Lemma~3$'$.
I claim only that, not self-refutation.

\paragraph{Objection 2 (van Fraassen's constructive empiricism).} ``Accept
that theories are truth-apt (strict instrumentalism is false) yet hold
that the aim of science is only empirical adequacy, not truth; this blocks
(T2) without operationalist or strict-instrumentalist commitments.''
\emph{Reply.} This concedes \S3's conclusive core and \S4.1--4.2; it
disputes only the abductive limb (T2), which I have already classified as
abductive, so it leaves the demonstrative results untouched. Against (T2)
specifically I do \emph{not} rely on the withdrawn slogan that there is no
theory-neutral observable. I rely on the no-miracles-to-structure argument
of \S5: empirical adequacy leaves the systematic novel-predictive and
unifying success of mature theories an unexplained coincidence, whereas
(structurally restricted) explanatory truth accounts for it and survives
the meta-induction. This lowers neither party to incoherence; it makes
(T2), in its structural form, the better explanation. I do not overstate:
van Fraassen's position is consistent, and my endorsement of (T2) is
correspondingly abductive, not conclusive.

\paragraph{Objection 3 (deflationism about truth).} ``If truth is merely
disquotational, the appeal to truth-aptness proves nothing metaphysical.''
\emph{Reply.} The conclusive results of this essay (\ref{C0}, Claims~4--5,
T1) require only the bipolarity/token-independence of contingent
description (Lemma~2), which deflationism preserves: ``$R$ is true'' is
intersubstitutable with $R$ and so inherits $R$'s token-independent
success condition. Deflationism is therefore no threat to the conclusive
core. The \emph{abductive} claims (semantic realism, T2) do trade on more
than disquotation---specifically on truth's explanatory role (\S3.3(ii),
\S5)---and I present them only as abductive; a thoroughgoing deflationist
who denies truth any explanatory role will resist them, and is, on the
present accounting, a coherent abductive opponent, not a refuted one.

\section{Conclusion}

The descriptive function of language is the use of declarative sentences
to say how things contingently stand, and it is the seat of truth, error,
and refutable inference (\S2). Built into that function is a standard of
correctness the act of describing cannot itself secure---\emph{token-independence}
(Lemma~2, \ref{C0})---and this is established conclusively. Token-independence
does not by itself entail recognition-transcendence; the stronger semantic
realism is established not analytically but as the best explanation of
descriptive and scientific practice, against a consistent sophisticated
anti-realist (\S3.3). The conclusive result alone refutes strict
operationalism and strict instrumentalism and establishes the
theory-ladenness of inquiry (\S4); the further, abductive result supports
the explanatory aim of theory, which the pessimistic meta-induction forces
us to take in \emph{structural} form (\S5). Hence the replacement of
operationalism and instrumentalism by theoreticism is justified:
deductively in its negative and theory-ladenness components, and by a
tempered inference to the best explanation in its explanatory-aim
component. The argument for realism from descriptive language is genuine
and, at the level of token-independence, conclusive; the route from there
to full theoreticism is partly demonstrative and partly---honestly
labelled---abductive.

\end{document}
